\chapter{Smooth links and stable minimal cones}\label{sec:stable-cones}
Cone stability combines a spherical operator with a one-dimensional Hardy inequality. Here $n\geq2$ is the cone dimension; its link has dimension $n-1$ and lies in the unit sphere $\mathbb S^n\subset\mathbb R^{n+1}$. We assume a smooth closed connected embedded two-sided link, so the cone is smooth away from its vertex. These hypotheses exclude more general cones with singular links. Sources are \cite{ChMin}, Section 26, pp.~64--68, and the Simons identities in \cite{HP99}, Section 2.

\begin{definition}[The cone over a spherical hypersurface]\label{def:smooth-link-cone}
For a link $L\subset\mathbb S^n$ as above, set
\[
 C^*(L)=\{rp:r>0,\ p\in L\},\qquad C(L)=C^*(L)\cup\{0\}.
\]
The normal $\nu_L$ to $L$ within the sphere extends to the cone by $\nu_C(rp)=\nu_L(p)$. The point $0$ is not included in the smooth domain $C^*(L)$. This is the cone parametrization used in \cite{ChMin}, Section 26.
\end{definition}

\begin{proposition}[Metric, curvature, and Laplacian of a cone]\label{prop:cone-geometry}
Under $F(r,p)=rp$,
\[
 g_C=dr^2+r^2g_L,\qquad dA_C=r^{n-1}dr\,dA_L,
 \qquad H_C=r^{-1}H_L,\qquad |A_C|^2=r^{-2}|A_L|^2,
\]
\[
 \Delta_C v=\partial_{rr}v+\frac{n-1}{r}\partial_rv+r^{-2}\Delta_Lv.
\]
In particular $C^*(L)$ is minimal precisely when $L$ is minimal in the unit sphere.
\end{proposition}
\begin{proof}
The radial vector is $p$ and an angular coordinate vector is $r\partial_ip$. They are orthogonal and have the displayed metric. The normal is independent of $r$, so the radial principal curvature is zero; angular principal curvatures are those of the spherical link divided by $r$. The Euclidean derivative of $\nu_L$ in a link direction has no radial component, because $\langle\nu_L,p\rangle=\langle\nu_L,\partial_ip\rangle=0$. This gives the mean-curvature and norm formulas. The determinant of the metric is $r^{2(n-1)}\det g_L$; the coordinate divergence formula for the Laplacian yields the last identity. Compare \cite{ChMin}, the calculation preceding Proposition 26.1.
\end{proof}

\begin{definition}[Stability away from a cone vertex]\label{def:cone-punctured-stability}
For a minimal link $L$, the cone is stable away from its vertex if
\[
 Q_C(\phi,\phi)=\int_{C^*(L)}
 \left(|\nabla_C\phi|^2-|A_C|^2\phi^2\right)dA_C\geq0
 \quad(\phi\in C_c^\infty(C^*(L))).
\]
Thus every test is supported in an annulus with positive inner radius. This definition does not assert regularity at the vertex or area minimization among integral-current or finite-perimeter competitors. It is the test class in \cite{ChMin}, Proposition 26.1.
\end{definition}

\begin{lemma}[The sharp radial Hardy inequality]\label{lem:radial-hardy}
Let $n\geq2$ and $a=(n-2)/2$. For $\eta\in C_c^\infty((0,\infty))$,
\[
 \int_0^\infty|\eta'|^2r^{n-1}dr
 \geq a^2\int_0^\infty\eta^2r^{n-3}dr.
\]
The constant $a^2$ is sharp, including the value zero when $n=2$, and equality holds only for $\eta=0$ in this compactly supported class.
\end{lemma}
\begin{proof}
Put $t=\log r$ and $\psi(t)=r^a\eta(r)$. Direct substitution gives
\[
 \int_0^\infty\bigl(|\eta'|^2r^{n-1}-a^2\eta^2r^{n-3}\bigr)dr
 =\int_{\mathbb R}|\psi'|^2dt,
 \qquad \int_0^\infty\eta^2r^{n-3}dr=\int_{\mathbb R}\psi^2dt.
\]
The cross term integrates to zero because $\psi$ is compactly supported. A zero right side makes $\psi$ constant and thus zero. To prove sharpness, fix a nonzero $\chi\in C_c^\infty(\mathbb R)$ and use $\psi_T(t)=\chi(t/T)$. Its derivative-energy to $L^2$-norm ratio is $T^{-2}\int|\chi'|^2/\int\chi^2\to0$. This is the logarithmic substitution underlying \cite{ChMin}, Proposition 26.1, pp.~64--65.
\end{proof}

\begin{definition}[The spherical link Rayleigh quotient]\label{def:link-rayleigh}
For a smooth minimal link define
\[
 \lambda_0(L)=\inf_{0\ne\psi\in C^\infty(L)}
 \frac{\int_L(|\nabla\psi|^2-|A_L|^2\psi^2)\,dA_L}
 {\int_L\psi^2\,dA_L}.
\]
This is the lowest Rayleigh value of $-\Delta_L-|A_L|^2$. It omits the ambient Ricci term $n-1$ that appears in the Jacobi operator for $L$ as a hypersurface of $\mathbb S^n$. Consequently it is not that spherical hypersurface's full stability eigenvalue. See \cite{ChMin}, Section 26, p.~64.
\end{definition}

\begin{proposition}[The spectral criterion for cone stability]\label{prop:cone-spectral-stability}
The cone $C^*(L)$ is stable in Definition \ref{def:cone-punctured-stability} if and only if
\[
 \lambda_0(L)\geq-\frac{(n-2)^2}{4}.
\]
This is \cite{ChMin}, Proposition 26.1, attributed there to Simons.
\end{proposition}
\begin{proof}
Write $a=(n-2)/2$, $t=\log r$, and $v(t,p)=r^a\phi(r,p)$. The cone metric and the same integration by parts as in Lemma \ref{lem:radial-hardy} give the exact identity
\[
 Q_C(\phi,\phi)=\int_{\mathbb R}\int_L
 \left(|\partial_tv|^2+|\nabla_Lv|^2+(a^2-|A_L|^2)v^2\right)dA_Ldt.
\]
If $\lambda_0\geq-a^2$, the angular part plus $a^2v^2$ has nonnegative integral at each fixed $t$, proving stability. Conversely, if a smooth $\psi$ has Rayleigh quotient below $-a^2$, choose $v(t,p)=\chi(t/T)\psi(p)$, with nonzero smooth compactly supported $\chi$. After division by $\int\chi(t/T)^2dt$, the positive time-derivative term is $T^{-2}\int|\chi'|^2/\int\chi^2$ times $\int_L\psi^2$. For large $T$ the total form is negative. This contradicts stability and proves necessity without assuming that the Rayleigh infimum is attained.
\end{proof}

\begin{example}[Product-sphere links and the dimension-seven threshold]\label{ex:clifford-cone-stability}
Let $p,q\geq1$, set $n=p+q+1$, and put
\[
 L=\mathbb S^p\left(\sqrt{\frac p{p+q}}\right)
 \times\mathbb S^q\left(\sqrt{\frac q{p+q}}\right)\subset\mathbb S^n.
\]
Writing its radii as $a,b$, a unit normal is $\nu(x,y)=(bx/a,-ay/b)$. Its principal curvatures are $b/a$ with multiplicity $p$ and $-a/b$ with multiplicity $q$. Hence $H_L=0$ and $|A_L|^2=p+q=n-1$. The Rayleigh quotient is minimized by constants, so $\lambda_0=-(n-1)$. Proposition \ref{prop:cone-spectral-stability} shows that its cone is stable away from zero exactly when
$n-1\leq(n-2)^2/4$, or $n\geq7$ in this family. This is a stability calculation, not a classification of which product cones minimize area. For equal factors with $p=q\geq3$, Theorem \ref{thm:simons-cone-minimizing} proves the stronger finite-perimeter minimizing property. Compare the cone examples in \cite{ChMin}, Sections 13, 17, and 26.
\end{example}

\begin{theorem}[Simons identity for a minimal spherical hypersurface]\label{thm:spherical-simons}
Let $L^d\subset\mathbb S^{d+1}$ be a smooth two-sided minimal hypersurface of the unit sphere. Its scalar second fundamental form satisfies
\[
 \Delta_L h=(d-|A|^2)h,\qquad
 \frac12\Delta_L|A|^2=|\nabla A|^2+d|A|^2-|A|^4.
\]
Here $\Delta_Lh$ is the rough covariant Laplacian. These are the constant-curvature, minimal specializations of \cite{HP99}, Corollary 2.2, p.~49; the corresponding norm inequality is \cite{ChMin}, Lemma 26.5.
\end{theorem}
\begin{proof}
The Gauss and Codazzi equations follow by decomposing the ambient covariant derivative into its tangential and normal parts. In the unit sphere they give $K_L(e_i,e_k)=1+\kappa_i\kappa_k$ in a principal frame and $\nabla_kh_{ij}=\nabla_ih_{kj}$. The latter holds because the normal component of the ambient curvature on three tangent vectors is zero. Tracing Codazzi and using $H=0$ gives $\sum_k\nabla_kh_{kj}=0$.

At a chosen point use a geodesic orthonormal frame diagonalizing $h$. Commutation on a covariant two-tensor gives
\[
 \Delta h_{ij}=\sum_k\bigl((\nabla_k\nabla_i-\nabla_i\nabla_k)h\bigr)_{kj},
\]
\[
 \bigl((\nabla_k\nabla_i-\nabla_i\nabla_k)h\bigr)(e_k,e_j)
 =-h(R^L(e_k,e_i)e_k,e_j)-h(e_k,R^L(e_k,e_i)e_j).
\]
The diagonal component is therefore
\[
 \Delta h_{ii}=\sum_k(1+\kappa_i\kappa_k)(\kappa_i-\kappa_k)
 =d\kappa_i-\kappa_i|A|^2,
\]
since $\sum_k\kappa_k=0$. The off-diagonal components are zero at this point by the same Gauss equation in a principal frame. This proves the tensor identity. Take its inner product with $h$ and apply the product rule to $|A|^2$ to obtain the second identity. The curvature sign is the one fixed in Proposition \ref{prop:riemannian-mean-curvature-variation}.
\end{proof}

\begin{proposition}[An upper bound for the nonflat link Rayleigh value]\label{prop:nonflat-link-rayleigh}
For a smooth closed minimal link $L^{n-1}$ with $A_L$ not identically zero,
\[
 \lambda_0(L)\leq-(n-1).
\]
\end{proposition}
\begin{proof}
Integrate the spherical Simons identity to get
\[
 \int_L|\nabla A|^2-\int_L|A|^4
 =-(n-1)\int_L|A|^2.
\]
The Lipschitz function $w=|A|$ satisfies $|\nabla w|\leq|\nabla A|$ almost everywhere. Use $w$ in the Rayleigh quotient by smooth approximation, for example $w_\varepsilon=\sqrt{|A|^2+\varepsilon^2}$ and passage in $W^{1,2}$. Its denominator tends to the positive number $\int|A|^2$, and its numerator is no larger in the limit than the left side above. This proves the bound. The approximation avoids dividing by $|A|$ at its zero set. Compare \cite{ChMin}, the proof of Corollary 26.6, pp.~67--68.
\end{proof}

\begin{theorem}[Stable cones with smooth links are flat in dimensions at most six]\label{thm:simons-smooth-link-cones}
Let $2\leq n\leq6$ and let $L^{n-1}\subset\mathbb S^n$ be a smooth closed connected embedded two-sided minimal hypersurface. If $C^*(L)$ is stable away from its vertex, then $L$ is a great sphere and $C(L)$ is a hyperplane. This is the smooth-link version of the dimension bound in \cite{ChMin}, Corollary 26.6.
\end{theorem}
\begin{proof}
If $A_L$ is not identically zero, the previous proposition and the stability criterion give
\[
 -\frac{(n-2)^2}{4}\leq\lambda_0(L)\leq-(n-1),
\]
which is impossible for $2\leq n\leq6$. Hence $A_L=0$. Its spherical unit normal has zero Euclidean derivative in every tangent direction, so it is constant on connected $L$. Thus $L$ lies in the equator orthogonal to that normal. It is open in that equator by equality of dimensions, and closed by compactness; hence it is the entire great sphere. The hypotheses include smoothness of the link. Extending this statement to cones with additional singularities requires a separate regularity or dimension-reduction argument.
\end{proof}
